% GENERATED FILE. Edit canonical lessons, then run npm run generate:books.
% canonical-source-hash: fnv1a64:9702e67cb7debcb6
% canonical-lessons: PT-C445-andaime, PT-C445-grua, PT-C445-estaleiro, PT-C445-desabar, PT-C445-demolir

\chapter{Obras na rua}
\label{ch:pt-obras-na-rua}
\hlchaptermodality{445}

\begin{chapteropening}
\textbf{By the end of this chapter:} \emph{I can talk about building work: the scaffolding, the crane, the building site, collapsing and knocking things down.}

Complete the last lesson of chapter 445: I can talk about building work: the scaffolding, the crane, the building site, collapsing and knocking things down.
\end{chapteropening}

\section[o andaime]{o andaime — the scaffolding}

\label{lesson:PT-C445-andaime}



\begin{quote}
Before the new one: say the Portuguese for the bollard, then the Portuguese for to obstruct.
\end{quote}

\subsection*{You'll want to know: o andaime}
\textbf{o andaime} — \textquotedblleft{}the scaffolding\textquotedblright{}.

\begin{quote}
Há andaimes à frente do prédio.
\end{quote}

\begin{quote}
\emph{There is scaffolding in front of the building.}
\end{quote}

Usually plural: \emph{os andaimes}.

\begin{cousinweb}
\emph{andaime} is usually traced to Arabic \emph{ad-daʿā'im} 'the supports', but the origin is not settled; EtymDB lists Spanish \emph{andamio} as a cousin.
\end{cousinweb}

\begin{grammarlens}[title={What you've built}]
One more word for telling, explaining and describing at B1.
\end{grammarlens}

\subsection*{Your turn}
\begin{itemize}
\raggedright
  \item \emph{Say it:} \emph{o andaime}
  \item \emph{Say it:} \emph{o andaime}, once more
  \item \emph{Say it:} a sentence of your own with \emph{o andaime}
  \item \emph{Recall:} say the Portuguese for the bollard, then the Portuguese for to obstruct, then say \emph{o andaime} again
\end{itemize}

\begin{tcolorbox}[breakable,colback=teal!4,colframe=teal!35!black,title={Before you move on}]
What does o andaime mean? (\textquotedblleft{}The scaffolding\textquotedblright{}.) Say it once more.
\end{tcolorbox}
\section[a grua]{a grua — the crane (for building)}

\label{lesson:PT-C445-grua}



\begin{quote}
Before the new one: say the Portuguese for to obstruct, then the Portuguese for the scaffolding.
\end{quote}

\subsection*{You'll want to know: a grua}
\textbf{a grua} — \textquotedblleft{}the crane (for building)\textquotedblright{}.

\begin{quote}
Uma grua enorme levanta os tijolos.
\end{quote}

\begin{quote}
\emph{A huge crane lifts the bricks.}
\end{quote}

The bird is \emph{o grou}; the machine took its name from its long neck.

\begin{cousinweb}
\emph{grua} goes back to Latin \emph{grūs} 'crane, the bird' (EtymDB): the machine has the bird's long neck.
\end{cousinweb}

\begin{grammarlens}[title={What you've built}]
One more word for telling, explaining and describing at B1.
\end{grammarlens}

\subsection*{Your turn}
\begin{itemize}
\raggedright
  \item \emph{Say it:} \emph{a grua}
  \item \emph{Say it:} \emph{a grua}, once more
  \item \emph{Say it:} a sentence of your own with \emph{a grua}
  \item \emph{Recall:} say the Portuguese for to obstruct, then the Portuguese for the scaffolding, then say \emph{a grua} again
\end{itemize}

\begin{tcolorbox}[breakable,colback=teal!4,colframe=teal!35!black,title={Before you move on}]
What does a grua mean? (\textquotedblleft{}The crane (for building)\textquotedblright{}.) Say it once more.
\end{tcolorbox}
\section[o estaleiro]{o estaleiro — the building site; the shipyard}

\label{lesson:PT-C445-estaleiro}



\begin{quote}
Before the new one: say the Portuguese for the scaffolding, then the Portuguese for the crane.
\end{quote}

\subsection*{You'll want to know: o estaleiro}
\textbf{o estaleiro} — \textquotedblleft{}the building site; the shipyard\textquotedblright{}.

\begin{quote}
O estaleiro fecha ao fim de semana.
\end{quote}

\begin{quote}
\emph{The building site is closed at the weekend.}
\end{quote}

\emph{Os estaleiros navais} — the shipyards.

\begin{cousinweb}
\emph{estaleiro} came from Spanish \emph{astillero} 'shipyard', which came from Old French \emph{astelier} 'workshop' (EtymDB).
\end{cousinweb}

\begin{grammarlens}[title={What you've built}]
One more word for telling, explaining and describing at B1.
\end{grammarlens}

\subsection*{Your turn}
\begin{itemize}
\raggedright
  \item \emph{Say it:} \emph{o estaleiro}
  \item \emph{Say it:} \emph{o estaleiro}, once more
  \item \emph{Say it:} a sentence of your own with \emph{o estaleiro}
  \item \emph{Recall:} say the Portuguese for the scaffolding, then the Portuguese for the crane, then say \emph{o estaleiro} again
\end{itemize}

\begin{tcolorbox}[breakable,colback=teal!4,colframe=teal!35!black,title={Before you move on}]
What does o estaleiro mean? (\textquotedblleft{}The building site; the shipyard\textquotedblright{}.) Say it once more.
\end{tcolorbox}
\section[desabar]{desabar — to collapse, to fall down}

\label{lesson:PT-C445-desabar}



\begin{quote}
Before the new one: say the Portuguese for the crane, then the Portuguese for the building site.
\end{quote}

\subsection*{You'll want to know: desabar}
\textbf{desabar} — \textquotedblleft{}to collapse, to fall down\textquotedblright{}.

\begin{quote}
Com a chuva, um muro velho desabou.
\end{quote}

\begin{quote}
\emph{In the rain, an old wall collapsed.}
\end{quote}

Also of a storm: \emph{desabou uma tempestade}.

\begin{cousinweb}
\emph{desabar} is usually explained as \emph{des-} and \emph{aba} 'brim, edge', but the origin is not settled; EtymDB has no entry for it.
\end{cousinweb}

\begin{grammarlens}[title={What you've built}]
One more word for telling, explaining and describing at B1.
\end{grammarlens}

\subsection*{Your turn}
\begin{itemize}
\raggedright
  \item \emph{Say it:} \emph{desabar}
  \item \emph{Say it:} \emph{desabar}, once more
  \item \emph{Say it:} a sentence of your own with \emph{desabar}
  \item \emph{Recall:} say the Portuguese for the crane, then the Portuguese for the building site, then say \emph{desabar} again
\end{itemize}

\begin{tcolorbox}[breakable,colback=teal!4,colframe=teal!35!black,title={Before you move on}]
What does desabar mean? (\textquotedblleft{}To collapse, to fall down\textquotedblright{}.) Say it once more.
\end{tcolorbox}
\section[demolir]{demolir — to knock down, to demolish}

\label{lesson:PT-C445-demolir}



\begin{quote}
Before the new one: say the Portuguese for the building site, then the Portuguese for to collapse.
\end{quote}

\subsection*{You'll want to know: demolir}
\textbf{demolir} — \textquotedblleft{}to knock down, to demolish\textquotedblright{}.

\begin{quote}
Vão demolir a fábrica antiga.
\end{quote}

\begin{quote}
\emph{They are going to knock down the old factory.}
\end{quote}

The noun is \emph{a demolição}.

\begin{cousinweb}
\emph{demolir} is a learned borrowing from Latin \emph{dēmōlīrī} 'to pull down' (Lewis \& Short), built on \emph{mōlēs} 'mass, structure'.
\end{cousinweb}

\begin{grammarlens}[title={What you've built}]
One more word for telling, explaining and describing at B1.
\end{grammarlens}

\subsection*{Your turn}
\begin{itemize}
\raggedright
  \item \emph{Say it:} \emph{demolir}
  \item \emph{Say it:} \emph{demolir}, once more
  \item \emph{Say it:} a sentence of your own with \emph{demolir}
  \item \emph{Recall:} say the Portuguese for the building site, then the Portuguese for to collapse, then say \emph{demolir} again
\end{itemize}

\begin{tcolorbox}[breakable,colback=teal!4,colframe=teal!35!black,title={Before you move on}]
What does demolir mean? (\textquotedblleft{}To knock down, to demolish\textquotedblright{}.) Say it once more.
\end{tcolorbox}
